\documentclass[10pt]{article}

\usepackage[letterpaper,margin=0.58in]{geometry}
\usepackage{graphicx}
\usepackage{caption}
\usepackage{subcaption}
\usepackage{amsmath}
\usepackage{amssymb}
\usepackage{booktabs}
\usepackage{enumitem}
\usepackage[colorlinks=true,linkcolor=blue,citecolor=blue,urlcolor=blue]{hyperref}

\captionsetup{font=footnotesize,labelfont=bf}
\captionsetup[subfigure]{font=scriptsize,skip=0.1em}
\setlength{\parskip}{0.28em}
\setlength{\parindent}{0pt}
\setlength{\textfloatsep}{0.45em}
\setlength{\floatsep}{0.35em}
\setlength{\intextsep}{0.35em}
\setlist{nosep,leftmargin=1.2em}

\title{\vspace{-2.2em}Pseudo-time Control and Trajectory-Inference Benchmarking}
\author{Langtian Ma}
\date{}
\begin{document}
\maketitle
\vspace{-2.0em}

This note follows the previous simulator summary by updating the pseudo-time control component. The base model remains the same: a disentangled VAE maps log-normalized expression into latent space, and a latent diffusion model samples new latents before VAE decoding. The updated runs use whitening-recoloring affine interpolation to generate trajectories with known ground-truth pseudo-time and topology, then use those generated datasets to benchmark trajectory-inference (TI) methods.

\section*{Pseudo-time Interpolation}

Let \(z=(z_{\mathrm{bio}}, z_{\mathrm{tech}})\), where \(z_{\mathrm{bio}}\) contains biological and residual latent dimensions and \(z_{\mathrm{tech}}\) contains technical blocks such as batch. Given start and end state latents with empirical moments \((\mu_0,\Sigma_0)\) and \((\mu_1,\Sigma_1)\), we define an affine endpoint map in the biological subspace,
\[
T(z)=\mu_1+B(z-\mu_0),
\qquad
B=\Sigma_1^{1/2}\Sigma_0^{-1/2},
\]
and generate a pseudo-time path by sample-linear interpolation,
\[
z_{\mathrm{bio}}(\alpha)=(1-\alpha)z_{\mathrm{bio}}+\alpha T(z_{\mathrm{bio}}),
\qquad \alpha\in[0,1].
\]
The full latent vector \(\tilde z_\alpha=(z_{\mathrm{bio}}(\alpha),z_{\mathrm{tech}})\) is decoded to expression space, and the known control value \(\alpha\) is recorded as ground-truth pseudo-time.

To simulate non-linear topologies such as bifurcations, we compose affine interpolation paths that share an initial segment. Given a root state \(A\), a shared waypoint \(W\), and two terminal fates \(B\) and \(C\), we first generate the trunk path \(A \to W\), then generate the two daughter paths \(W \to B\) and \(W \to C\). This construction gives a known branch point and assigns every generated cell both a ground-truth pseudo-time and a lineage label.

Beyond producing a bifurcation with known topology, we want a continuous knob for how different the two daughter branches look, from nearly indistinguishable (low discrepancy) to strongly diverging (high discrepancy). This lets us stress-test TI methods on their ability to resolve branches, not just order cells within one branch.

\paragraph{1. Direction and length \((u,r)\): geometric knob.} For the two daughter branches \(W \to B\) and \(W \to C\), each branch is characterized by a displacement vector
\[
d_B = r_B u_B, \qquad d_C = r_C u_C,
\]
from the shared start \(W\) to each endpoint. Here \(u_B\) and \(u_C\) are unit vectors giving the two branch directions, while \(r_B \ge 0\) and \(r_C \ge 0\) are the corresponding branch lengths. For example, if \(B\) and \(C\) are observed endpoint states, one can set \(u_B=(\mu_B-\mu_W)/\|\mu_B-\mu_W\|\), \(r_B=\|\mu_B-\mu_W\|\), and analogously for \(u_C\) and \(r_C\). Given a start distribution with moments \((\mu_W, \Sigma_W)\) and endpoint covariances \(\Sigma_B\) and \(\Sigma_C\), the target endpoint moments are:

\[
\mu_{\text{target},B} = \mu_W + d_B, \qquad \Sigma_B,
\qquad
\mu_{\text{target},C} = \mu_W + d_C, \qquad \Sigma_C.
\]

Each branch is generated by applying the affine interpolation primitive from \((\mu_W, \Sigma_W)\) to its corresponding endpoint moments, with \(\alpha \in [0, 1]\) indexing pseudo-time along that branch. The cosine similarity \(\langle u_B, u_C \rangle\) then measures how similar the two daughter directions are. The lengths \(r_B\) and \(r_C\) independently control how far the two branches travel. We can control the branch discrepancy and length by varying \(\langle u_B, u_C \rangle\) and \((r_B, r_C)\)

\paragraph{2. Branch-point position \(\tau\): topological knob.} We can pick an observed state \(W\) to act as a shared waypoint and split at pseudo-time \(\tau \in [0, 1]\). Run three independent affine interpolations, \(A \to W\), \(W \to B\), and \(W \to C\), each over \(\alpha \in [0, 1]\), and map them onto a common pseudo-time axis by rescaling:

\[
t_{\text{trunk}} = \tau \cdot \alpha, \qquad t_{\text{branch}} = \tau + (1 - \tau) \cdot \alpha.
\]

This rescaling is purely a relabelling of pseudo-time and leaves the per-segment interpolation unchanged. Cell counts are allocated proportionally (trunk \(\propto \tau\), each branch \(\propto 1 - \tau\)) to keep density roughly uniform in pseudo-time. \(\tau \to 0\) recovers a split at the root; \(\tau \to 1\) gives a near-linear trajectory with a very late split.

\paragraph{3. Noise scale \(\sigma\): SNR knob.} To vary branch separability without changing the trajectory geometry, we add per-branch isotropic Gaussian noise with scale \(\sigma\). Larger \(\sigma\) lowers the signal-to-noise ratio and makes the same underlying branches harder to recover.

\textbf{Interpolation results.} The branch-direction knob controls daughter-branch discrepancy: lower cosine-simlarity makes the two terminal arms more separated in UMAP space (Figure~\ref{fig:pt-control}a,b). A dose-response check over \(\alpha\in[0,1]\) shows increasing separation between generated pseudo-time groups (Figure~\ref{fig:pt-control}c). For branching, changing \(\tau\) visibly moves the split along the common pseudo-time axis: larger \(\tau\) creates a denser trunk and sparser branches (Figure~\ref{fig:pt-control}d).

\begin{figure}[t]
  \centering
  \begin{minipage}[c]{0.40\linewidth}
    \begin{subfigure}[t]{\linewidth}
      \centering
      \includegraphics[width=\linewidth]{../../../experiments/outputs/2026-04-21/15-25-00_branch_direction_knob/results/w_1.00_umap.png}
      \caption{Direction discrepancy, \(\langle u_B, u_C \rangle = 0.789\)}
    \end{subfigure}
    \vspace{0.15em}

    \begin{subfigure}[t]{\linewidth}
      \centering
      \includegraphics[width=\linewidth]{../../../experiments/outputs/2026-04-21/15-25-00_branch_direction_knob/results/w_2.00_umap.png}
      \caption{Direction discrepancy, \(\langle u_B, u_C \rangle = 0.246\)}
    \end{subfigure}
  \end{minipage}
  \hspace{0.02\linewidth}
  \begin{minipage}[c]{0.40\linewidth}
    \begin{subfigure}[t]{\linewidth}
      \centering
      \includegraphics[width=\linewidth]{figures/pt_dose_response_curve.png}
      \caption{Pseudo-time dose response}
    \end{subfigure}
  \end{minipage}

  \vspace{0.25em}

  \begin{subfigure}[t]{0.9\linewidth}
    \centering
    \includegraphics[width=\linewidth]{figures/tau_comparison_umap.png}
    \caption{Branch point \(\tau\)}
  \end{subfigure}

  \caption{Pseudo-time control outputs. Panels (a) and (b) show branch-direction discrepancy control at \(\langle u_B, u_C \rangle = 0.789\) and \(\langle u_B, u_C \rangle = 0.246\). Panel (c) quantifies increasing pseudo-time separation as \(\alpha\) moves from start to endpoint. Panel (d) shows controllable branch-point placement in a Ductal-to-Ngn3-high-EP-to-Alpha/Beta bifurcation.}
  \label{fig:pt-control}
\end{figure}

\section*{TI Methods Benchmarking}

We benchmark Scanpy DPT/PAGA, Slingshot, and Monocle3 on generated trajectories with known pseudo-time, lineage, and branch structure. The sweeps vary three controlled difficulty axes: branch discrepancy, branch-point position \(\tau\), and noise scale \(\sigma\). Performance is summarized by global Spearman correlation for pseudo-time ordering and adjusted Rand index (ARI) for lineage recovery (Figure~\ref{fig:ti-benchmark}).

\textbf{Endpoint discrepancy sweep.} As the branches separate, ordering generally becomes easier, especially for DPT/PAGA and Monocle3. Lineage recovery improves more weakly and remains close to unresolved for several methods.

\textbf{Branch-point sweep.} DPT/PAGA is the most stable across early and intermediate branch points, while Slingshot becomes competitive for late splits but does not recover lineage labels reliably.

\textbf{Noise sweep.} Increasing \(\sigma\) lowers the branch signal-to-noise ratio and produces the clearest difficulty ladder. DPT/PAGA performs best under low-to-moderate noise, but all methods degrade as noise becomes large.

\begin{figure}[htbp]
  \centering
  \begin{subfigure}[t]{0.28\linewidth}
    \centering
    \includegraphics[width=\linewidth]{figures/ti_discrepancy_curves.png}
    \caption{Branch discrepancy}
  \end{subfigure}
  \hfill
  \begin{subfigure}[t]{0.28\linewidth}
    \centering
    \includegraphics[width=\linewidth]{figures/ti_tau_curves.png}
    \caption{Branch point \(\tau\)}
  \end{subfigure}
  \hfill
  \begin{subfigure}[t]{0.28\linewidth}
    \centering
    \includegraphics[width=\linewidth]{figures/ti_noise_curves.png}
    \caption{Noise scale \(\sigma\)}
  \end{subfigure}
  \caption{TI benchmark results across three difficulty axes. Each panel reports global pseudo-time Spearman correlation and lineage ARI for Scanpy DPT/PAGA, Slingshot, and Monocle3.}
  \label{fig:ti-benchmark}
\end{figure}

\end{document}
